%%%%%%%%%%%%%%%%
%%% exod 14 %%%
%%%%%%%%%%%%%%%%

\Chapter{14}

\Verse{1}
{
	\hP{וַיְדַבֵּר יְהֹוָה אֶל מֹשֶׁה לֵּאמֹר׃}
}
{
	\eP{And YHWH spoke to Moses, saying,}
}
{
}

\Verse{2}
{
	\hP{דַּבֵּר אֶל בְּנֵי יִשְׂרָאֵל וְיָשֻׁבוּ וְיַחֲנוּ לִפְנֵי פִּי הַחִירֹת בֵּין מִגְדֹּל וּבֵין הַיָּם לִפְנֵי בַּעַל צְפֹן נִכְחוֹ תַחֲנוּ עַל הַיָּם׃}
}
{
	\eP{
		“Speak to the children of Israel that they should go back
		and camp in front of Pi-Hahiroth, between Migdol and the
		sea, in front of Baal-Zephon. You shall camp facing it, by
		the sea.
	}
}
{
}

\Verse{3}
{
	\hP{וְאָמַר פַּרְעֹה לִבְנֵי יִשְׂרָאֵל נְבֻכִים הֵם בָּאָרֶץ סָגַר עֲלֵיהֶם הַמִּדְבָּר׃}
}
{
	\eP{
		And Pharaoh will say about the children of Israel:
		‘They’re muddled in the land! The wilderness has closed
		them in.’
	}
}
{
}

\Verse{4}
{
	\hP{וְחִזַּקְתִּי אֶת לֵב פַּרְעֹה וְרָדַף אַחֲרֵיהֶם וְאִכָּבְדָה בְּפַרְעֹה וּבְכל חֵילוֹ וְיָדְעוּ מִצְרַיִם כִּי אֲנִי יְהֹוָה וַיַּעֲשׂוּ כֵן׃}
}
{
	\eP{
		And I’ll strengthen Pharaoh’s heart, and he’ll pursue
		them, and I’ll be glorified against Pharaoh and against
		all of his army, and Egypt will know that I am YHWH.” And
		they did so.
	}
}
{
}

\Verse{5}
{
	\hJ{וַיֻּגַּד לְמֶלֶךְ מִצְרַיִם כִּי בָרַח הָעָם וַיֵּהָפֵךְ לְבַב פַּרְעֹה וַעֲבָדָיו אֶל הָעָם וַיֹּאמְרוּ מַה זֹּאת עָשִׂינוּ כִּי שִׁלַּחְנוּ אֶת יִשְׂרָאֵל מֵעבְדֵנוּ׃}
}
{
	\eJ{
		And it was told to the king of Egypt that the people had
		fled, and the heart of Pharaoh and his servants was
		changed toward the people. And they said, “What is this
		that we’ve done, that we let Israel go from serving us?!”
	}
}
{
%	\footnote{
%		and the heart of Pharaoh … was changed. The word that was used for the
%		stick’s changing into a snake and the water’s changing into blood and
%		the changing of the wind that carried away the locusts is now used to
%		describe Pharaoh’s change of heart. Those past cases, plus the passive
%		(Niphal) form of the verb here, indicate that it is not Pharaoh himself
%		but God who causes the change.
%	}
}

\Verse{6}
{
	\hJ{וַיֶּאְסֹר אֶת רִכְבּוֹ וְאֶת עַמּוֹ לָקַח עִמּוֹ׃}
}
{
	\eJ{And he hitched his chariot and took his people with him.}
}
{
}

\Verse{7}
{
	\hJ{וַיִּקַּח שֵׁשׁ מֵאוֹת רֶכֶב בָּחוּר וְכֹל רֶכֶב מִצְרָיִם וְשָׁלִשִׁם עַל כֻּלּוֹ׃}
}
{
	\eJ{
		And he took six hundred chosen chariots—and all the
		chariotry of Egypt—and officers over all of it.
	}
}
{
}

\Verse{8}
{
	\hP{וַיְחַזֵּק יְהֹוָה אֶת לֵב פַּרְעֹה מֶלֶךְ מִצְרַיִם וַיִּרְדֹּף אַחֲרֵי בְּנֵי יִשְׂרָאֵל וּבְנֵי יִשְׂרָאֵל יֹצְאִים בְּיָד רָמָה׃}
}
{
	\eP{
		And YHWH strengthened the heart of Pharaoh, king of Egypt,
		and he pursued the children of Israel. And the children of
		Israel were going out with a high hand.
	}
}
{
}

\Verse{9}
{
	\hP{וַיִּרְדְּפוּ מִצְרַיִם אַחֲרֵיהֶם וַיַּשִּׂיגוּ אוֹתָם חֹנִים עַל הַיָּם כּל סוּס רֶכֶב פַּרְעֹה וּפָרָשָׁיו וְחֵילוֹ עַל פִּי הַחִירֹת לִפְנֵי בַּעַל צְפֹן׃}
}
{
	\eP{
		And Egypt pursued them, and they caught up to them,
		camping by the sea—every chariot horse of Pharaoh and his
		horsemen and his army—at Pi-Hahiroth, in front of
		Baal-Zephon.
	}
}
{
}

\Verse{10}
{
	\hJ{וּפַרְעֹה הִקְרִיב וַיִּשְׂאוּ בְנֵי יִשְׂרָאֵל אֶת עֵינֵיהֶם וְהִנֵּה מִצְרַיִם נֹסֵעַ אַחֲרֵיהֶם וַיִּירְאוּ מְאֹד וַיִּצְעֲקוּ בְנֵי יִשְׂרָאֵל אֶל יְהֹוָה׃}
}
{
	\eJ{
		And Pharaoh came close! And the children of Israel raised
		their eyes, and here was Egypt coming after them, and they
		were very afraid. And the children of Israel cried out to
		YHWH.
	}
}
{
}

\Verse{11}
{
	\hJ{וַיֹּאמְרוּ אֶל מֹשֶׁה הֲמִבְּלִי אֵין קְבָרִים בְּמִצְרַיִם לְקַחְתָּנוּ לָמוּת בַּמִּדְבָּר מַה זֹּאת עָשִׂיתָ לָּנוּ לְהוֹצִיאָנוּ מִמִּצְרָיִם׃}
}
{
	\eJ{
		And they said to Moses, “Was it because of an
		absence—none!—of graves in Egypt that you took us to die
		in the wilderness?! What is this that you’ve done to us to
		bring us out of Egypt?
	}
}
{
%	\footnote{
%		an absence—none!—of graves. Their words are redundant. The prophet
%		Elijah uses the same expression: “Is it because of an absence—none!—of a
%		God in Israel that you go to inquire of Baal?!” (2 Kings 1:3,6,16). The
%		redundancy dramatizes the question in each case, as if to say, “Was it
%		because of a complete and utter lack …”—the way people speak when
%		accusing someone of doing something completely senseless. Footnote
%		14:11. you’ve done to us. As with Joseph, people focus on the human,
%		Moses, as if he and not God were the one who brought this about.
%	}
}

\Verse{12}
{
	\hJ{הֲלֹא זֶה הַדָּבָר אֲשֶׁר דִּבַּרְנוּ אֵלֶיךָ בְמִצְרַיִם לֵאמֹר חֲדַל מִמֶּנּוּ וְנַעַבְדָה אֶת מִצְרָיִם כִּי טוֹב לָנוּ עֲבֹד אֶת מִצְרַיִם מִמֻּתֵנוּ בַּמִּדְבָּר׃}
}
{
	\eJ{
		Isn’t this the thing that we spoke to you in Egypt,
		saying: Stop from us! And let’s serve Egypt. Because
		serving Egypt is better for us than our dying in the
		wilderness!”
	}
}
{
}

\Verse{13}
{
	\hJ{וַיֹּאמֶר מֹשֶׁה אֶל הָעָם אַל תִּירָאוּ הִתְיַצְּבוּ וּרְאוּ אֶת יְשׁוּעַת יְהֹוָה אֲשֶׁר יַעֲשֶׂה לָכֶם הַיּוֹם כִּי אֲשֶׁר רְאִיתֶם אֶת מִצְרַיִם הַיּוֹם לֹא תֹסִפוּ לִרְאֹתָם עוֹד עַד עוֹלָם׃}
}
{
	\eJ{
		And Moses said to the people, “Don’t be afraid. Stand
		still and see YHWH’s salvation that He’ll do for you
		today. For, as you’ve seen Egypt today, you’ll never see
		them again, ever.
	}
}
{
}

\Verse{14}
{
	\hJ{יְהֹוָה יִלָּחֵם לָכֶם וְאַתֶּם תַּחֲרִשׁוּן׃}
}
{
	\eJ{YHWH will fight for you, and you’ll keep quiet!”}
}
{
}

\Verse{15}
{
	\hP{וַיֹּאמֶר יְהֹוָה אֶל מֹשֶׁה מַה תִּצְעַק אֵלָי דַּבֵּר אֶל בְּנֵי יִשְׂרָאֵל וְיִסָּעוּ׃}
}
{
	\eP{
		And YHWH said to Moses, “Why do you cry out to me? Speak
		to the children of Israel that they should move!
	}
}
{
}

\Verse{16}
{
	\hP{וְאַתָּה הָרֵם אֶת מַטְּךָ וּנְטֵה אֶת יָדְךָ עַל הַיָּם וּבְקָעֵהוּ וְיָבֹאוּ בְנֵי יִשְׂרָאֵל בְּתוֹךְ הַיָּם בַּיַּבָּשָׁה׃}
}
{
	\eP{
		And you, lift your staff and reach your hand out over the
		sea—and split it! And the children of Israel will come
		through the sea on the dry ground.
	}
}
{
}

\Verse{17}
{
	\hP{וַאֲנִי הִנְנִי מְחַזֵּק אֶת לֵב מִצְרַיִם וְיָבֹאוּ אַחֲרֵיהֶם וְאִכָּבְדָה בְּפַרְעֹה וּבְכל חֵילוֹ בְּרִכְבּוֹ וּבְפָרָשָׁיו׃}
}
{
	\eP{
		And I, here, I’m strengthening Egypt’s heart, and they’ll
		come after them, and I’ll be glorified against Pharaoh and
		against all of his army, against his chariots and against
		his horsemen.
	}
}
{
}

\Verse{18}
{
	\hP{וְיָדְעוּ מִצְרַיִם כִּי אֲנִי יְהֹוָה בְּהִכָּבְדִי בְּפַרְעֹה בְּרִכְבּוֹ וּבְפָרָשָׁיו׃}
}
{
	\eP{
		And Egypt will know that I am YHWH when I’m glorified
		against Pharaoh, against his chariots and against his
		horsemen.”
	}
}
{
}

\Verse{19}
{
	\hE{וַיִּסַּע מַלְאַךְ הָאֱלֹהִים הַהֹלֵךְ לִפְנֵי מַחֲנֵה יִשְׂרָאֵל וַיֵּלֶךְ מֵאַחֲרֵיהֶם וַיִּסַּע עַמּוּד הֶעָנָן מִפְּנֵיהֶם וַיַּעֲמֹד מֵאַחֲרֵיהֶם׃}
}
{
	\eE{
		And the angel of {\God} who was going in front of the camp of
		Israel moved and went behind them, and the column of cloud
		went from in front of them and stood behind them.
	}
}
{
}

\Verse{20}
{
	\hE{וַיָּבֹא בֵּין מַחֲנֵה מִצְרַיִם וּבֵין מַחֲנֵה יִשְׂרָאֵל וַיְהִי הֶעָנָן וְהַחֹשֶׁךְ וַיָּאֶר אֶת הַלָּיְלָה וְלֹא קָרַב זֶה אֶל זֶה כּל הַלָּיְלָה׃}
}
{
	\eE{
		And it came between the camp of Egypt and the camp of
		Israel. And there was the cloud and darkness [for the
		Egyptians], while it [the column of fire] lit the night
		{}[for the Israelites], and one did not come near the other
		all night.
	}
}
{
}

\Verse{21}
{
	\hJ{וַיֵּט מֹשֶׁה אֶת יָדוֹ עַל הַיָּם וַיּוֹלֶךְ יְהֹוָה אֶת הַיָּם בְּרוּחַ קָדִים עַזָּה כּל הַלַּיְלָה וַיָּשֶׂם אֶת הַיָּם לֶחָרָבָה וַיִּבָּקְעוּ הַמָּיִם׃}
}
{
	\eJ{
		And Moses reached his hand out over the sea, and YHWH
		drove back the sea with a strong east wind all night and
		turned the sea into dry ground, and the sea was split.
	}
}
{
}

\Verse{22}
{
	\hP{וַיָּבֹאוּ בְנֵי יִשְׂרָאֵל בְּתוֹךְ הַיָּם בַּיַּבָּשָׁה וְהַמַּיִם לָהֶם חוֹמָה מִימִינָם וּמִשְּׂמֹאלָם׃}
}
{
	\eP{
		And the children of Israel came through the sea on the dry
		ground. And the water was a wall to them at their right
		and at their left.
	}
}
{
}

\Verse{23}
{
	\hP{וַיִּרְדְּפוּ מִצְרַיִם וַיָּבֹאוּ אַחֲרֵיהֶם כֹּל סוּס פַּרְעֹה רִכְבּוֹ וּפָרָשָׁיו אֶל תּוֹךְ הַיָּם׃}
}
{
	\eP{
		And Egypt pursued and came after them, every horse of
		Pharaoh, his chariots and his horsemen, through the sea.
	}
}
{
}

\Verse{24}
{
	\hJ{וַיְהִי בְּאַשְׁמֹרֶת הַבֹּקֶר וַיַּשְׁקֵף יְהֹוָה אֶל מַחֲנֵה מִצְרַיִם בְּעַמּוּד אֵשׁ וְעָנָן וַיָּהם אֵת מַחֲנֵה מִצְרָיִם׃}
}
{
	\eJ{
		And it was in the morning watch, and YHWH gazed at Egypt’s
		camp through a column of fire and cloud and threw Egypt’s
		camp into tumult,
	}
}
{
}

\Verse{25}
{
	\hJ{וַיָּסַר אֵת אֹפַן מַרְכְּבֹתָיו וַיְנַהֲגֵהוּ בִּכְבֵדֻת וַיֹּאמֶר מִצְרַיִם אָנוּסָה מִפְּנֵי יִשְׂרָאֵל כִּי יְהֹוָה נִלְחָם לָהֶם בְּמִצְרָיִם׃}
}
{
	\eJ{
		and He turned its chariots’ wheel so that it drove it with
		heaviness. And Egypt said, “Let me flee from Israel,
		because YHWH is fighting for them against Egypt!”
	}
}
{
%	\footnote{
%		heaviness. The recurrence of the word “heavy” to describe yet another
%		event conveys that the plagues section of the story and the Red Sea
%		section are united. The event at the sea is not just one more episode;
%		it is a climax to what has been happening.
%	}
}

\Verse{26}
{
	\hP{וַיֹּאמֶר יְהֹוָה אֶל מֹשֶׁה נְטֵה אֶת יָדְךָ עַל הַיָּם וְיָשֻׁבוּ הַמַּיִם עַל מִצְרַיִם עַל רִכְבּוֹ וְעַל פָּרָשָׁיו׃}
}
{
	\eP{
		And YHWH said to Moses, “Reach your hand out over the sea,
		and the water will go back over Egypt, over his chariots
		and over his horsemen.”
	}
}
{
}

\Verse{27}
{
	\hP{וַיֵּט מֹשֶׁה אֶת יָדוֹ עַל הַיָּם וַיָּשׁב הַיָּם לִפְנוֹת בֹּקֶר לְאֵיתָנוֹ וּמִצְרַיִם נָסִים לִקְרָאתוֹ וַיְנַעֵר יְהֹוָה אֶת מִצְרַיִם בְּתוֹךְ הַיָּם׃}
}
{
	\eP{
		And Moses reached his hand out over the sea, and the sea
		went back to its strong flow toward morning, and Egypt was
		fleeing toward it. And YHWH tossed the Egyptians into the
		sea.
	}
}
{
}

\Verse{28}
{
	\hP{וַיָּשֻׁבוּ הַמַּיִם וַיְכַסּוּ אֶת הָרֶכֶב וְאֶת הַפָּרָשִׁים לְכֹל חֵיל פַּרְעֹה הַבָּאִים אַחֲרֵיהֶם בַּיָּם לֹא נִשְׁאַר בָּהֶם עַד אֶחָד׃}
}
{
	\eP{
		And the waters went back and covered the chariots and the
		horsemen—all of Pharaoh’s army who were coming after them
		in the sea. Not even one of them was left.
	}
}
{
}

\Verse{29}
{
	\hP{וּבְנֵי יִשְׂרָאֵל הָלְכוּ בַיַּבָּשָׁה בְּתוֹךְ הַיָּם וְהַמַּיִם לָהֶם חֹמָה מִימִינָם וּמִשְּׂמֹאלָם׃}
}
{
	\eP{
		And the children of Israel had gone on the dry ground
		through the sea, and the water had been a wall to them at
		their right and at their left.
	}
}
{
}

\Verse{30}
{
	\hR{וַיּוֹשַׁע יְהֹוָה בַּיּוֹם הַהוּא אֶת יִשְׂרָאֵל מִיַּד מִצְרָיִם וַיַּרְא יִשְׂרָאֵל אֶת מִצְרַיִם מֵת עַל שְׂפַת הַיָּם׃}
}
{
	\eR{
		And YHWH saved Israel from Egypt’s hand that day. And
		Israel saw Egypt dead on the seashore,
	}
}
{
}

\Verse{31}
{
	\hR{וַיַּרְא יִשְׂרָאֵל אֶת הַיָּד הַגְּדֹלָה אֲשֶׁר עָשָׂה יְהֹוָה בְּמִצְרַיִם וַיִּירְאוּ הָעָם אֶת יְהֹוָה וַיַּאֲמִינוּ בַּיהֹוָה וּבְמֹשֶׁה עַבְדּוֹ׃}
}
{
	\eR{
		and Israel saw the big hand that YHWH had used against
		Egypt, and the people feared YHWH, and they trusted in
		YHWH and in Moses His servant.
	}
}
{
%	\footnote{
%		they trusted. This word is also translated often as “believed” (Hebrew
%		wayya’mînû). It does not have the meaning here that it has in later
%		religious concepts. That is, it does not function in the sense of
%		believing that a God exists. This notion of belief in does not occur in
%		Biblical Hebrew (nor in other ancient Near Eastern languages). In pagan
%		religion the gods, being observable forces in nature (e.g., the sun, the
%		sky, the storm wind), are not a matter of belief but of knowledge. So in
%		the conception of YHWH in Exodus, God becomes known; God’s existence and
%		power are a matter of knowledge, not belief. When one has seen ten
%		plagues and a sea split and has a column of cloud and fire visible at
%		all times, one does not ask, “Do you believe in God?” As the term is
%		used in the Hebrew Bible, it means not belief in, but belief that; that
%		is, it means that if YHWH says He will do something one can trust that
%		He will do it.
%	}
}
